\scanpage{109}
\margin{\mspage{N27}}
Thus the subalgebra $\Lambda_{1,n}$ of the model precisely determines the nilpotent quotient $\pi_1/C_n$ tensored by $Q$ or $R$.

It is perhaps important to note that the minimal model $\Lambda$ of the de Rham algebra is constructed in a purely algebraic computational construction as indicated in the text above. The structure of the model then determines the form of the lower central series of $\pi_1$ (over $Q$ and $R$) and so on for the further rational or real homotopy theory.

The subalgebra $\Lambda_1$ of $\Lambda$ mapping to the de Rham algebra is constructed so we have an isomorphism of $H^1$ and an injection on $H^2$. Then 2 dim'l generators are added in a possibly infinite sequence of steps. The first step creates an isomorphism on $H^2$ and the later steps produce an injection on $H^3$. Then 3 dim'l generators are added in a possibly infinite sequence of steps, 4 dim'l, and so on to build the complete model.

The structure of model over $Q$ for the rational de Rham algebra of a space corresponds (via the Hirsch calculation described in Appendix H) to a tower of spaces and maps. Each map in the tower is a principal fibration where the fibre has one non-zero homotopy group which is a rational vector space.

This tower of spaces is the rational homotopy type of the space which is defined in general for all spaces regardless of the fundamental group.

We have described above the relationship between the $\Lambda_1$ of the model (and therefore the tower) and the lower central series of the fundamental group of the space. The significance of the higher dimensional structure of the model and the tower for the general non-simply connected spaces is somewhat complicated and not yet completely understood except in a philosophical way.
\margin{108}
